% QPU React/TypeScript Parity Specification
% Brings react-qpu-mvp up to feature parity with BPForbes/QPU (Python)
% Compile: pdflatex qpu-react-python-parity.tex (twice for TOC)

\documentclass[11pt,a4paper]{article}

\usepackage[margin=1in]{geometry}
\usepackage{amsmath,amssymb,amsfonts,mathtools}
\usepackage{booktabs}
\usepackage{enumitem}
\usepackage{hyperref}
\usepackage{listings}
\usepackage{xcolor}
\usepackage{graphicx}
\usepackage{tikz}
\usetikzlibrary{arrows.meta,positioning,shapes.geometric}

\hypersetup{
  colorlinks=true,
  linkcolor=blue!60!black,
  citecolor=green!50!black,
  urlcolor=blue!70!black,
}

\lstdefinestyle{protocol}{
  basicstyle=\ttfamily\small,
  breaklines=true,
  frame=single,
  backgroundcolor=\color{gray!8},
}
\lstdefinestyle{code}{
  basicstyle=\ttfamily\footnotesize,
  breaklines=true,
  frame=single,
  backgroundcolor=\color{gray!5},
  language=Java,
}

\title{\textbf{QPU React Parity Specification}\\[0.4em]
\large Physics Engine, Mathematics, API, Circuit Builder, Hilbert Space, and Data Migration}
\author{react-qpu-mvp $\rightarrow$ BPForbes/QPU alignment}
\date{\today}

\newcommand{\C}{\mathbb{C}}
\newcommand{\R}{\mathbb{R}}
\newcommand{\ket}[1]{\left|#1\right\rangle}
\newcommand{\bra}[1]{\left\langle#1\right|}
\newcommand{\ketbra}[2]{\ket{#1}\!\bra{#2}}
\newcommand{\Tr}{\operatorname{Tr}}
\newcommand{\HS}{\mathcal{H}_{\mathrm{log}}}
\newcommand{\xor}{\oplus}

\begin{document}
\maketitle

\begin{abstract}
The React/TypeScript QPU application (\texttt{react-qpu-mvp}) is deployed as a \emph{static site} (Vite $\rightarrow$ GitHub Pages). Simulation, compilation, and correction currently execute entirely in the browser. The reference Python implementation (\texttt{BPForbes/QPU}) provides cycle-oriented AST evaluation, a register-indexed \texttt{HilbertSpace} audit log, hybrid sparse/dense state vectors, noise channels, and checkpointing (\texttt{SAVE\_STATE}/\texttt{LOAD\_STATE}). This document specifies the mathematical, architectural, API, UI, and data-object upgrades required to bring the React stack to parity while preserving static-site deployment via an optional sidecar API and in-browser fallback. \textbf{Verified against current code} (2026): \texttt{engine.ts}, \texttt{particleTracking.ts}, \texttt{blochQuadrature.ts}, \texttt{qpuAst.ts}, \texttt{qpucirFile.ts}; HilbertSpace and sidecar are \emph{specified}, not yet implemented.
\end{abstract}

\tableofcontents
\newpage

%=============================================================================
\section{Scope and gap analysis}
%=============================================================================

\subsection{Current React capabilities (verified)}

\begin{itemize}[noitemsep]
  \item Full state-vector simulation over $\ket{\psi}\in\C^{2^n}$ as \texttt{Complex[]} (\texttt{engine.ts}).
  \item Compile-time lowering of \texttt{.qpucir} to flat \texttt{CircuitGate[]} (\texttt{qpuAst.ts}); \texttt{INCREASECYCLE} flushes workspace zeros at compile time, not at runtime.
  \item Marginal Bloch analytics: $\rho_k$ via \texttt{marginalRho}, $\vec{B}_k$, purity $r=\|\vec{B}_k\|$, $\Tr(\rho_k^2)=(1+r^2)/2$, triple-integral $\langle\rho\rangle$ on the Bloch ball (\texttt{particleTracking.ts}, \texttt{blochQuadrature.ts}).
  \item Ket display: $\ket{\psi}=\cos(\theta/2)\ket{0}+e^{i\phi}\sin(\theta/2)\ket{1}$ via \texttt{ketFromSpherical} (Euler form for $\alpha,\beta$).
  \item Truth-table testing via \texttt{.qpuio} and \texttt{moduleTestApi.ts}.
  \item Correction Lab: truth-table-driven synthesis and LLM intent parsing.
  \item Export payloads: \texttt{qpucir} v1 JSON (\texttt{createQpucirPayload}), \texttt{qpuio} v1 JSON, catalog v1 in \texttt{localStorage}.
  \item \texttt{SAVE\_STATE}/\texttt{LOAD\_STATE} parsed in \texttt{qpuAst.ts} but \textbf{not lowered} to runtime (no-op at compile time).
\end{itemize}

\subsection{Python capabilities not yet realized in React}

\begin{itemize}[noitemsep]
  \item Runtime \textbf{HilbertSpace} logger: $\HS\colon (\mathit{register}\times\mathbb{N})\to\C^{2^k}$.
  \item Cycle-by-cycle AST interpreter with \texttt{is\_ready} dependency scheduling.
  \item Per-register \texttt{local\_states} synchronized with global \texttt{StateVector}.
  \item Hybrid \textbf{sparse/dense} state representation with copy-on-write generations.
  \item \textbf{Noise}: amplitude damping, phase damping, depolarizing, readout error; idle noise per cycle.
  \item \textbf{Device fingerprint}: per-qubit $f$, $T_1$, $T_2$ property table.
  \item Runtime \textbf{checkpoints} (\texttt{SAVE\_STATE}/\texttt{LOAD\_STATE}) and runtime fingerprints.
  \item Ket formatting for arbitrary register dimension (\texttt{format\_qubit\_state}).
\end{itemize}

\subsection{Parity principle}

\begin{quote}
\textbf{P1 (semantic parity):} For every supported protocol command, the observable outputs (measurements, return values, truth-table rows, and Hilbert log entries) must match Python within tolerance $\varepsilon_a=10^{-6}$ on amplitudes and $\varepsilon_p=10^{-4}$ on probabilities.

\textbf{P2 (static-site preservation):} The GitHub Pages bundle remains a client; parity features that require shared state or heavy jobs are exposed through an optional HTTP API and degrade gracefully to in-browser execution when no API URL is configured.

\textbf{P3 (forward-compatible data):} All export formats gain an explicit \texttt{version} field and migration path; protected bundled \texttt{.qpuio} truth tables are never mutated in place.
\end{quote}

%=============================================================================
\section{Physics engine upgrades}
%=============================================================================

\subsection{Dual execution modes}

Define two engine backends selectable at compile or session creation:

\begin{enumerate}[label=\textbf{E\arabic*.}]
  \item \textbf{Flat mode} (current): \texttt{compileQpuProtocol} $\rightarrow$ \texttt{runCircuit}. Fast, UI-friendly; Hilbert log is reconstructed post hoc from gate list.
  \item \textbf{Cycle mode} (Python parity): AST nodes remain scheduled; \texttt{INCREASECYCLE} advances clock $c\mapsto c+1$; gates evaluate only when memory dependencies at cycle $c$ are satisfied.
\end{enumerate}

Let $c\in\mathbb{N}$ be the simulation cycle and $M_c$ the memory map (register $\rightarrow$ state at cycle $c$). Cycle mode satisfies:
\begin{equation}
  \text{node.evaluate}(M_c, c, \mathrm{QPU}, \HS) \quad\text{iff}\quad \text{node.is\_ready}(M_c, c).
\end{equation}

\subsection{Global state and register views}

Python maintains:
\begin{align}
  \ket{\Psi}_{\mathrm{global}} &\in \C^{2^{n_{\mathrm{phys}}}}, \\
  \ket{\psi_r} &\in \C^{2^{k_r}}, \quad r\in\{\text{physical qubits}\}\cup\{\texttt{0p},\texttt{1p},\texttt{sp},\ldots\}.
\end{align}

\textbf{Upgrade:} Introduce \texttt{RegisterBank}:
\begin{itemize}
  \item \texttt{globalState}: $\C^{2^n}$ (authoritative for entangling gates on physical wires).
  \item \texttt{localStates}: map $r\mapsto\C^{2}$ for physical indices.
  \item \texttt{customStates}: map token $\mapsto\C^{2^k}$.
  \item \texttt{rebuildGlobal()} via Kronecker product consistent with Python \texttt{rebuild\_global\_state}.
  \item \texttt{extractQubit(k)}: partial trace marginal for wire $k$, normalized.
\end{itemize}

After every single-qubit gate on physical wire $k$, refresh $\ket{\psi_k}\leftarrow\texttt{extractQubit}(k)$.

\subsection{Sparse/dense hybrid (optional Phase~2)}

Let $S=\{i : |a_i|>\tau\}$ with $\tau=10^{-12}$. Use sparse storage when $|S|\le \eta\cdot 2^n$ with $\eta=0.05$ (Python \texttt{DENSE\_THRESHOLD}). Copy-on-write generation counter avoids cloning on read-only simulation branches (correction search).

\subsection{Noise model}

When \texttt{noise.enabled}, apply after each \texttt{INCREASECYCLE} (mirroring \texttt{qpu\_base.py}):

\textbf{Amplitude damping} with $p_1 = 1 - e^{-\Delta t / T_1}$:
\begin{align}
  E_0 &= \begin{bmatrix}1&0\\0&\sqrt{1-p_1}\end{bmatrix}, &
  E_1 &= \begin{bmatrix}0&\sqrt{p_1}\\0&0\end{bmatrix}.
\end{align}

\textbf{Phase damping} with $p_\phi = 1 - e^{-\Delta t / T_2}$:
\begin{align}
  K_0 &= \sqrt{1-p_\phi}\, I, &
  K_1 &= \sqrt{p_\phi}\begin{bmatrix}1&0\\0&-1\end{bmatrix}.
\end{align}
Apply stochastically: choose $K_0$ or $K_1$ with probability $\|K_j\ket{\psi}\|^2$, then renormalize.

\textbf{Depolarizing} (per idle step, probability $p_{\mathrm{dep}}$):
\begin{equation}
  \rho \mapsto (1-p_{\mathrm{dep}})\,\rho + \frac{p_{\mathrm{dep}}}{2}\, I.
\end{equation}
Equivalently: with probability $p_{\mathrm{dep}}$, apply random Pauli $X$, $Y$, or $Z$ (as in Python \texttt{apply\_idle\_noise}).

\textbf{Readout error} on \texttt{MEASURE}: classical bit-flip with probability $p_{\mathrm{ro}}$ (flip reported outcome $0\leftrightarrow 1$ while leaving collapsed state as measured, matching Python \texttt{readout\_error}).

\textbf{Device table} $\mathbf{P}\in\R^{n\times 3}$: columns $(f_g, T_1, T_2)$ per qubit, serializable in session exports.

\subsection{Checkpointing}

For checkpoint name $s$, store:
\begin{equation}
  \mathcal{C}(s) = \bigl(\ket{\Psi},\, M_c,\, c,\, \mathrm{measurements},\, \HS\big|_{\le c}\bigr).
\end{equation}

\texttt{SAVE\_STATE} $s$ writes $\mathcal{C}(s)$; \texttt{LOAD\_STATE} $s$ restores. Compiler must \emph{lower} these opcodes to runtime hooks in cycle mode (today accepted but no-ops in \texttt{qpuAst.ts} lines 513--514).

\subsection{Particle tracking integration}

Retain existing Bloch pipeline when \texttt{trackParticles: true} (\texttt{engine.ts}). Extend snapshots with Hilbert cycle index:
\begin{equation}
  \mathrm{Snapshot}(k,c) = \bigl(\vec{B}_k,\, \ket{\psi_k^{\mathrm{marg}}},\, \rho_k,\, c\bigr).
\end{equation}

%=============================================================================
\section{Mathematical foundations and handling}
%=============================================================================

\subsection{Number system}

All amplitudes use IEEE-754 \texttt{float64}. Define:
\begin{align}
  \texttt{Complex} &= \{z = x + iy : x,y\in\R\}, \\
  \|z\|^2 &= x^2 + y^2, \quad |z| = \sqrt{\|z\|^2}.
\end{align}

\textbf{Normalization:} After measurement collapse and noise Kraus maps, renormalize:
$\ket{\psi}\leftarrow \ket{\psi}/\|\psi\|$ if $\|\psi\|>\varepsilon_{\mathrm{nrm}}$ with $\varepsilon_{\mathrm{nrm}}=10^{-12}$.

\textbf{Tolerances:}
\begin{center}
\begin{tabular}{@{}ll@{}}
\toprule
Quantity & Tolerance \\
\midrule
Amplitude equality & $\varepsilon_a = 10^{-6}$ \\
Probability / purity display & $\varepsilon_p = 10^{-4}$ \\
Zero amplitude pruning & $\tau = 10^{-12}$ \\
Pure-state Bloch radius & $1 \pm 10^{-6}$ \\
\bottomrule
\end{tabular}
\end{center}

\textbf{Composition across $N$ gates:} For deterministic amplitude comparisons in regression tests, adopt the \textbf{worst-case linear bound} $|a-a'| \le N\cdot\varepsilon_a$ when comparing end-to-end state vectors after $N$ serialized gates. For probability/purity display checks, use $N\cdot\varepsilon_p$. If modeling independent per-gate numerical noise, an optional RSS bound $\sqrt{N}\cdot\varepsilon_a$ may be used for diagnostics but is \textbf{not} the default pass/fail rule (avoids spurious failures on long circuits). Pruning threshold $\tau$ and pure-state Bloch radius tolerance apply \textbf{per snapshot}, not scaled by $N$. Example: a 16-gate regression comparing $|\Psi\rangle$ uses amplitude tolerance $16\times 10^{-6}$; a single-gate Bloch radius check still uses $1\pm 10^{-6}$.

\subsection{Basis indexing convention}

Qubit $q$ (0-indexed, MSB-left in labels) occupies bit $(n-1-q)$ of basis index $i$:
\begin{equation}
  \mathrm{bit}(i,q,n) = \left\lfloor \frac{i}{2^{n-1-q}} \right\rfloor \bmod 2.
\end{equation}
This matches \texttt{hasBit} in \texttt{operations.ts}. \textbf{All API and export schemas must document MSB-left basis strings} (e.g.\ \texttt{"101"} for $n=3$).

\subsection{Single-qubit gates}

For gate matrix $U\in\C^{2\times 2}$ on wire $k$, let $m=2^{n-1-k}$. For each index $i$ with $\mathrm{bit}(i,k,n)=0$ and paired index $j=i\xor m$:
\begin{equation}
  a'_i = U_{00}\,a_i + U_{01}\,a_j, \qquad
  a'_j = U_{10}\,a_i + U_{11}\,a_j.
\end{equation}
Here $\xor$ denotes bitwise XOR (flip bit $k$). This matches \texttt{applySingleQubitGate} in \texttt{operations.ts}.

\subsection{Marginal and Bloch coordinates}

For qubit $k$, reduced density matrix (as implemented in \texttt{marginalRho}):
\begin{equation}
  \rho_k = \begin{bmatrix}
    \rho_{00} & \rho_{01} \\
    \rho_{01}^* & \rho_{11}
  \end{bmatrix}, \quad
  \rho_{ab} = \sum_{i:\,\mathrm{bit}(i,k,n)=a} a_i \overline{a_{i'}},
\end{equation}
where $i' = i \xor 2^{n-1-k}$ (same bit-flip pairing as single-qubit gates).

Bloch vector (Pauli expectations; matches \texttt{blochVectorForQubit} and \texttt{MATRIX\_Y} in \texttt{matrices.ts}):
\begin{equation}
  \vec{B}_k = \bigl(\Tr(\rho_k X),\, \Tr(\rho_k Y),\, \Tr(\rho_k Z)\bigr)
  = \bigl(2\Re\rho_{01},\, 2\Im\rho_{01},\, \rho_{00}-\rho_{11}\bigr).
\end{equation}

Spherical coordinates on the Bloch ball:
\begin{align}
  x &= r\sin\theta\cos\phi, &
  y &= r\sin\theta\sin\phi, &
  z &= r\cos\theta, &
  r &= \|\vec{B}_k\|_2.
\end{align}

Pure-state ket display (when $r\approx 1$; \texttt{ketFromSpherical}):
\begin{equation}
  \ket{\psi_k} = \cos\frac{\theta}{2}\ket{0} + e^{i\phi}\sin\frac{\theta}{2}\ket{1},
  \quad \alpha=\cos\frac{\theta}{2},\quad \beta=e^{i\phi}\sin\frac{\theta}{2}.
\end{equation}

\textbf{Entanglement caveat} (must appear in API docs and UI): $\vec{B}_k$ from $\rho_k$ describes the \emph{marginal} state; for entangled $\ket{\Psi}$, $\rho_k$ may be mixed even when the global state is pure. Today \texttt{snapshotParticle} still calls \texttt{ketFromSpherical} on $(\theta,\phi)$ even when $r<1$; extension: branch to eigenket display when mixed.

\subsection{Extending existing particle mathematics}

\subsubsection{Triple integral for particle density}

Current \texttt{blochBallRhoExpectationFast} computes
\begin{equation}
  \langle\rho\rangle = \frac{1}{V}\iiint_{\mathrm{ball}} \rho(r,\theta,\phi)\, r^2\sin\theta\, dr\,d\theta\,d\phi,\quad V=\frac{4\pi}{3},
\end{equation}
via separable Gauss--Legendre quadrature (10 radial $\times$ 10$\times$10 angular samples). Noise $\eta=1-r$ blends peaked kernel with uniform $\rho_{\mathrm{uni}}=3/(4\pi)$.

\textbf{Extensions} (same quadrature skeleton):
\begin{itemize}
  \item \textbf{Depolarizing kernel:} $\rho_{\mathrm{kernel}}=(1-p_{\mathrm{dep}})\delta_{\vec{r}_0}+p_{\mathrm{dep}}\rho_{\mathrm{uni}}$ tied to \texttt{NoiseConfig.depolarizingProb}.
  \item \textbf{Cycle diffusion:} spread $\sigma(c)$ grows with Hilbert cycle $c$ for time-smoothed inspector.
  \item \textbf{Exact mode:} optional voxel grid (as in \texttt{blochQuadrature.test.ts}) when high precision requested.
\end{itemize}

\subsubsection{Ket / Euler layer}

Extend \texttt{PsiKet}:
\begin{itemize}
  \item \textbf{Pure} ($r\ge 1-\varepsilon_a$): read $(\alpha,\beta)$ from amplitudes or $\rho_k$ with global phase fixed ($\alpha\in\R_{\ge 0}$).
  \item \textbf{Mixed} ($r<1-\varepsilon_a$): expose dominant eigenket of $\rho_k$ with weight $\lambda_1$; keep \texttt{formatPsiKet} for pure branch only.
\end{itemize}

\subsubsection{Gate-induced geometry}

Extend \texttt{particleDelta} with rotation angle $\omega$ where $\cos(\omega/2)=1-\|\Delta\vec{B}\|^2/2$ for $r\approx 1$, and fidelity $F(\rho,\rho')$ between snapshots for checkpoint validation.

\subsection{Ket formatting (Python parity)}

Implement \texttt{formatQubitState}($\vec{a}\in\C^d$):
\begin{itemize}
  \item $d=2$: use \texttt{formatPsiKet} / \texttt{formatComplex}.
  \item $d=2^n$: sum over indices with $|a_i|>\tau$: \texttt{"$a_k\ket{b_{n-1}\ldots b_0}$"}.
\end{itemize}

\subsection{Math module structure}

Refactor into layers (no UI imports):

\begin{center}
\begin{tabular}{@{}ll@{}}
\toprule
Module & Responsibility \\
\midrule
\texttt{math/complex.ts} & $\C$ arithmetic (existing) \\
\texttt{math/linalg.ts} & Kronecker, partial trace, normalization \\
\texttt{math/ketFormat.ts} & \texttt{formatQubitState}, basis labels \\
\texttt{physics/stateVector.ts} & dense (and optional sparse) backend \\
\texttt{physics/registerBank.ts} & per-register + global sync \\
\texttt{physics/hilbertSpace.ts} & $\HS$ logger \\
\texttt{physics/noise.ts} & Kraus channels, device table \\
\texttt{physics/particleTracking.ts} & Bloch (existing, extended) \\
\bottomrule
\end{tabular}
\end{center}

%=============================================================================
\section{Hilbert Space subsystem}
%=============================================================================

\subsection{Definition}

Python's \texttt{HilbertSpace} (\texttt{qpu/hilbert.py}) is \textbf{not} the Hilbert space $\C^{2^n}$ itself; it is an append-only \emph{observation log}:
\begin{equation}
  \HS \colon (\mathit{Reg} \times \mathbb{N}) \rightharpoonup \C^{2^k}
  \cup \{(\texttt{\_\_all\_\_}, c) \mapsto \ket{\Psi}_{\mathrm{global}}\}.
\end{equation}

\subsection{TypeScript interface}

\begin{lstlisting}[style=protocol]
type HilbertKey = { register: string | number; cycle: number };
type HilbertEntry = {
  key: HilbertKey;
  dimension: number;
  amplitudes: Complex[];
  ketFormatted: string;
  recordedAt: 'gate' | 'measure' | 'set' | 'checkpoint' | 'cycle';
};
type HilbertSpace = {
  output(register, cycle, state): void;
  outputAll(cycle, globalState): void;
  input(register, cycle): Complex[] | undefined;
  prune(maxCycle): void;
  view(): string;
  entries(): HilbertEntry[];
};
\end{lstlisting}

\subsection{Semantics of \texttt{input(register, cycle)}}

\textbf{Destructive pop} (Python parity): \texttt{input(r,c)} returns a copy of the \texttt{Complex[]} stored under \texttt{HilbertKey\{register:r, cycle:c\}} and \textbf{removes} that entry from internal storage. Subsequent calls with the same $(r,c)$ return \texttt{undefined}. Absent keys return \texttt{undefined} without mutation.

\textbf{Non-destructive reads} use \texttt{entries()} or a future \texttt{peek(register, cycle)} if needed; \texttt{output}/\texttt{outputAll} are append/overwrite-at-key writes; \texttt{prune(maxCycle)} deletes all entries with $\mathrm{cycle}>\mathrm{maxCycle}$.

Parity tests must assert pop semantics for \texttt{input} and append semantics for \texttt{output}.

\subsection{Reference implementation sketch}

\begin{lstlisting}[style=code]
// physics/hilbertSpace.ts (proposed)
export class HilbertSpace {
  private space = new Map<string, Complex[]>();

  private key(register: string | number, cycle: number): string {
    return `${String(register)}@${cycle}`;
  }

  output(register: string | number, cycle: number, state: Complex[]): void {
    this.space.set(this.key(register, cycle), state.map((z) => ({ ...z })));
  }

  outputAll(cycle: number, globalState: Complex[]): void {
    this.output('__all__', cycle, globalState);
  }

  input(register: string | number, cycle: number): Complex[] | undefined {
    const k = this.key(register, cycle);
    const value = this.space.get(k);
    if (!value) return undefined;
    this.space.delete(k);
    return value.map((z) => ({ ...z }));
  }

  prune(maxCycle: number): void {
    for (const k of [...this.space.keys()]) {
      const cycle = Number(k.split('@').pop());
      if (cycle > maxCycle) this.space.delete(k);
    }
  }

  entries(): HilbertEntry[] { /* map space -> HilbertEntry + ketFormatted */ }
}
\end{lstlisting}

\subsection{Recording policy}

Record to $\HS$ when:
\begin{enumerate}
  \item A gate mutates register $r$ at cycle $c$.
  \item \texttt{SET} assigns a definite state.
  \item \texttt{MEASURE} collapses $r$.
  \item \texttt{INCREASECYCLE} optionally calls \texttt{outputAll}.
  \item Child process \texttt{RUNCHILD} returns updated tokens.
\end{enumerate}

\subsection{UI and correction usage}

\begin{itemize}
  \item \textbf{Inspector panel}: sort by $(c, r)$; filter by register; diff cycles $c$ and $c+1$.
  \item \textbf{Correction}: when a truth-table row fails, attach minimal $\HS$ slice for implicated registers across cycles where outputs are read.
  \item \textbf{LLM context}: include \texttt{hilbert.view()} text capped at $N$ entries with guidance to prefer gate-level fixes over table edits for protected processes.
\end{itemize}

%=============================================================================
\section{State vector implementation}
%=============================================================================

\subsection{Current code (verified)}

\texttt{engine.ts} uses dense \texttt{Complex[]} length $2^n$, \texttt{createInitialState}, \texttt{applyGate} $\rightarrow$ \texttt{gates/registry}, optional \texttt{trackParticles} $\rightarrow$ \texttt{snapshotAllParticles}.

\subsection{Proposed \texttt{StateVector} wrapper}

\begin{lstlisting}[style=code]
// physics/stateVector.ts (proposed; mirrors Python statevector.py)
export class StateVector {
  private data: Complex[];
  private gen = 0;

  constructor(readonly qubitCount: number) {
    this.data = Array.from({ length: 2 ** qubitCount }, () => ZERO);
    this.data[0] = ONE;
  }

  clone(): StateVector {
    const next = new StateVector(this.qubitCount);
    next.data = this.data.map((z) => ({ ...z }));
    next.gen = this.gen;
    return next;
  }

  ensureCow(currentGen: number): void {
    if (this.gen !== currentGen) { this.data = this.data.map((z) => ({ ...z })); this.gen = currentGen; }
  }

  applySingleQubitGate(matrix: Complex[][], k: number): void {
    const n = this.qubitCount;
    const mask = 1 << (n - 1 - k);
    for (let i = 0; i < this.data.length; i++) {
      if (i & mask) continue;
      const j = i | mask;
      const a = this.data[i]; const b = this.data[j];
      this.data[i] = add(scale(a, matrix[0][0].re), scale(b, matrix[0][1].re)); // use complex mul
      this.data[j] = add(scale(a, matrix[1][0].re), scale(b, matrix[1][1].re));
    }
  }

  fullState(): Complex[] { return this.data.map((z) => ({ ...z })); }
}
\end{lstlisting}

Sparse/dense hybrid remains Phase~II; initial parity uses dense arrays identical to today's \texttt{engine.ts} loops.

%=============================================================================
\section{API changes}
%=============================================================================

\subsection{In-process API (\texttt{moduleTestApi.ts})}

Extend types:

\begin{lstlisting}[style=protocol]
type ExecutionMode = 'flat' | 'cycle';
type NoiseConfig = { enabled: boolean; depolarizingProb?: number;
  readoutError?: number; dt?: number };
type ModuleTestRequest = {
  // ... existing fields ...
  executionMode?: ExecutionMode;
  noise?: NoiseConfig;
  recordHilbert?: boolean;
  hilbertMaxCycle?: number;
};
type ModuleTestResponse = {
  // ... existing fields ...
  hilbert?: HilbertEntry[];
  checkpoints?: Record<string, CheckpointSummary>;
  deviceTable?: number[][];
};
\end{lstlisting}

\textbf{Correction with Hilbert Space:}
\begin{enumerate}
  \item Run failing test with \texttt{recordHilbert: true}, \texttt{executionMode: 'cycle'}.
  \item Build \texttt{CorrectionContext}: failing rows; $\HS$ entries for \texttt{RETURNVALS} at final cycle; \texttt{OperationTransition[]} when \texttt{trackParticles}.
  \item \texttt{correctCircuit} ranks fixes by truth-table pass rate first, then minimal $\sum_k \|\Delta\vec{B}_k\|$.
  \item Never emit intents that rewrite protected \texttt{.qpuio} tables.
\end{enumerate}

\subsection{Optional HTTP API (static-site sidecar)}

Environment: \texttt{VITE\_QPU\_API\_URL}. When unset, all calls stay in-process (current behavior).

\begin{center}
\begin{tabular}{@{}lll@{}}
\toprule
Method & Path & Purpose \\
\midrule
POST & \texttt{/v1/compile} & protocol $\rightarrow$ gates or AST \\
POST & \texttt{/v1/sessions} & create session (mode, noise, $n$) \\
DELETE & \texttt{/v1/sessions/\{id\}} & destroy session \\
POST & \texttt{/v1/sessions/\{id\}/run} & flat run \\
POST & \texttt{/v1/sessions/\{id\}/step} & one gate or one cycle \\
GET  & \texttt{/v1/sessions/\{id\}/hilbert} & $\HS$ entries \\
GET  & \texttt{/v1/sessions/\{id\}/bloch} & particle snapshots \\
POST & \texttt{/v1/sessions/\{id\}/checkpoint} & \texttt{SAVE\_STATE} \\
POST & \texttt{/v1/sessions/\{id\}/restore} & \texttt{LOAD\_STATE} \\
POST & \texttt{/v1/test} & \texttt{runModuleTest} \\
POST & \texttt{/v1/jobs} & async correction \\
\bottomrule
\end{tabular}
\end{center}

\subsubsection{Error handling and wire format}

\textbf{HTTP status codes:}
\begin{itemize}
  \item \texttt{200} success; \texttt{201} session/job created.
  \item \texttt{400} malformed JSON or schema validation failure.
  \item \texttt{404} unknown \texttt{sessionId} or checkpoint name.
  \item \texttt{409} session mutation conflict (see concurrency).
  \item \texttt{422} compilation error, invalid protocol, or unsupported \texttt{execution.mode}.
  \item \texttt{500} internal error; \texttt{503} sidecar at session capacity.
\end{itemize}

\textbf{Uniform error envelope} (all endpoints):
\begin{lstlisting}[style=protocol]
{ "error": {
    "code": "SESSION_NOT_FOUND",
    "message": "No session with id ...",
    "details": { ... optional structured fields ... }
  }}
\end{lstlisting}
Examples: \texttt{COMPILE\_ERROR}, \texttt{STATE\_CORRUPT}, \texttt{VERSION\_UNSUPPORTED}, \texttt{SESSION\_CONFLICT}.

\textbf{API versioning:} \texttt{/v1/} prefix is stable for minor additive fields. Breaking changes require \texttt{/v2/}. Clients send optional \texttt{Accept-Version: 1} header; servers include \texttt{X-QPU-API-Version: 1}. Deprecation: \texttt{Sunset} header + 6-month overlap documented in release notes.

\textbf{Amplitude wire format:}
\begin{itemize}
  \item JSON array of \texttt{\{"re": number, "im": number\}} objects; IEEE-754 \texttt{float64} semantics.
  \item JSON is text; \textbf{no binary endianness} on the wire. If a future binary codec is added, use big-endian IEEE-754 per field.
  \item \texttt{NaN}, \texttt{+Inf}, \texttt{-Inf} are \textbf{rejected} on input (\texttt{400}); never emitted on output (normalize or error).
  \item Length must be exact power of two $2^n$; mismatch $\rightarrow$ \texttt{422}.
\end{itemize}
Ket strings may be precomputed server-side for bandwidth; amplitudes remain authoritative.

\subsubsection{Session concurrency and lifecycle}

\textbf{Single-writer per session:} \texttt{POST .../run}, \texttt{POST .../step}, \texttt{POST .../checkpoint}, \texttt{POST .../restore} are serialized through a server-side mutex (FIFO queue). Concurrent mutating requests receive \texttt{409 SESSION\_CONFLICT} with \texttt{details.expectedRevision}.

\textbf{Optimistic revision:} each session has monotonic \texttt{revision} (integer). Mutating requests should include \texttt{If-Match: <revision>}; mismatch $\rightarrow$ \texttt{409}.

\textbf{Concurrent reads allowed:} \texttt{GET .../hilbert}, \texttt{GET .../bloch} may run parallel to each other but see last-committed revision (read-your-writes after mutations complete).

\textbf{Browser tabs:} default \textbf{one session per tab} (client generates \texttt{sessionId} in \texttt{sessionStorage}). Sharing a session across tabs is opt-in via URL hash; clients must coordinate writes (leader election or single-writer tab) to avoid \texttt{409} storms.

\textbf{Lifecycle:}
\begin{itemize}
  \item Idle timeout: 30 minutes (configurable); timer resets on any request.
  \item Explicit \texttt{DELETE /v1/sessions/\{id\}}; idempotent.
  \item Server GC: sweep expired sessions every 5 minutes.
  \item Default max concurrent sessions per sidecar instance: 64 (configurable); excess $\rightarrow$ \texttt{503}.
\end{itemize}

\subsection{Sidecar implementation sketch}

\begin{lstlisting}[style=code]
// server/index.ts (proposed Node sidecar)
import { createServer } from 'node:http';
import { compileQpuProtocol } from '../src/simulator/compiler/qpuAst';
import { runCircuit, stepCircuitGate } from '../src/simulator/engine';
import { runModuleTest } from '../src/simulator/moduleTestApi';
import { HilbertSpace } from '../src/simulator/physics/hilbertSpace';

type Session = {
  id: string; revision: number; lock: Promise<void>;
  qubitCount: number; state: Complex[]; measurements: MeasurementMap;
  hilbert: HilbertSpace; mode: 'flat' | 'cycle'; expiresAt: number;
};

const sessions = new Map<string, Session>();

function json(res, status, body) {
  res.writeHead(status, { 'Content-Type': 'application/json', 'X-QPU-API-Version': '1' });
  res.end(JSON.stringify(body));
}

createServer(async (req, res) => {
  if (req.url === '/v1/compile' && req.method === 'POST') {
    try {
      const { source, librarySources } = await readJson(req);
      const compiled = compileQpuProtocol(source, librarySources ?? {});
      json(res, 200, { compiled });
    } catch (e) { json(res, 422, { error: { code: 'COMPILE_ERROR', message: String(e) } }); }
    return;
  }
  // POST /v1/sessions, serialized POST .../step with revision check, etc.
}).listen(process.env.PORT ?? 8787);
\end{lstlisting}

Deploy sidecar separately from GitHub Pages; SPA reads \texttt{import.meta.env.VITE\_QPU\_API\_URL}.

%=============================================================================
\section{Circuit builder updates}
%=============================================================================

\subsection{Cycle-aware canvas}

\begin{itemize}
  \item Optional \textbf{cycle lanes} when \texttt{executionMode=cycle}.
  \item \texttt{INCREASECYCLE} as lane separators (compiler-internal zeros remain hidden).
  \item Token syntax \texttt{\$Name:c} from \texttt{tokenMap}.
\end{itemize}

\subsection{New palette items}

\begin{center}
\begin{tabular}{@{}lll@{}}
\toprule
Item & Visual & Runtime \\
\midrule
\texttt{SAVE\_STATE} & bookmark icon & checkpoint \\
\texttt{LOAD\_STATE} & restore icon & restore \\
\texttt{INCREASECYCLE} & dashed separator & $c\gets c+1$ \\
Noise toggle & toolbar & $\mathbf{P}$ editor \\
Hilbert inspector & side panel & live $\HS$ feed \\
\bottomrule
\end{tabular}
\end{center}

\subsection{Particle / Bloch view}

\begin{itemize}
  \item Scrubber binds to gate \texttt{step} or cycle $c$.
  \item Show $\Tr(\rho_k^2)$ and mixed badge when $r<1-\varepsilon_p$.
  \item Optional WebGL Bloch sphere; CSS fallback retained (\texttt{ParticleView.tsx}).
\end{itemize}

%=============================================================================
\section{Integration and usage}
%=============================================================================

\subsection{Boot sequence}

\begin{enumerate}
  \item Load catalog (\texttt{processCatalog.ts}).
  \item Compile with \texttt{executionMode} from settings.
  \item Create \texttt{SimulationSession}: \texttt{RegisterBank} + \texttt{StateVector} + optional \texttt{HilbertSpace} + \texttt{NoiseEngine}.
  \item Run / step $\rightarrow$ \texttt{ExecutionResult} + \texttt{hilbert.entries()}.
\end{enumerate}

\subsection{Feature flags}

\begin{lstlisting}[style=protocol]
type QpuFeatureFlags = {
  executionMode: 'flat' | 'cycle';
  recordHilbert: boolean;
  trackParticles: boolean;
  noise: boolean;
  remoteApi: boolean;
};
\end{lstlisting}

Defaults: \texttt{flat}, \texttt{recordHilbert=false}, \texttt{trackParticles=true}, \texttt{noise=false}.

%=============================================================================
\section{Data objects and migration}
\label{sec:exports}
%=============================================================================

\subsection{\texttt{qpucir} payload}

\textbf{v1} (current, \texttt{createQpucirPayload}):
\begin{lstlisting}[style=protocol]
{ format: "qpucir", version: 1, name, source,
  compiled: { qubitCount, gates, tokenMap }, exportedAt }
\end{lstlisting}

\textbf{v2} (proposed):
\begin{lstlisting}[style=protocol]
{
  format: "qpucir", version: 2, name, source, exportedAt,
  execution: { mode: "flat"|"cycle", noise?: NoiseConfig },
  compiled: { qubitCount, gates, tokenMap, processParams?, returnValues? },
  hilbertPolicy: "none" | "on-export" | "full-trace",
  hilbert?: HilbertEntry[],
  checkpoints?: Record<string, CheckpointPayload>,
  deviceTable?: number[][]
}
\end{lstlisting}

\subsection{Backward compatibility and schema validation}

\textbf{v1 readers encountering v2 files:}
\begin{itemize}
  \item \textbf{(A)} If \texttt{version===2} and runtime is v1-only: \textbf{reject} with error ``\texttt{qpucir v2 requires engine $\geq$2.x; upgrade client or re-export as v1}'' unless \texttt{execution.mode==="flat"} and \texttt{hilbertPolicy==="none"} and no \texttt{checkpoints}/\texttt{deviceTable}---then \textbf{strip unknown fields} and recompile from \texttt{source} only.
  \item \textbf{(B)} If \texttt{execution.mode==="cycle"} (or \texttt{hilbertPolicy!=="none"}) and runtime cannot execute those features: \textbf{fail fast} with descriptive error; do not silently downgrade.
\end{itemize}

\textbf{v2 validation rules} (on import):
\begin{enumerate}
  \item \texttt{format==="qpucir"} and \texttt{version===2}.
  \item \texttt{execution.mode} $\in$ \{\texttt{flat},\texttt{cycle}\}.
  \item \texttt{hilbertPolicy} $\in$ \{\texttt{none},\texttt{on-export},\texttt{full-trace}\}.
  \item If \texttt{hilbertPolicy==="none"} then \texttt{hilbert} must be absent.
  \item If \texttt{hilbertPolicy!=="none"} then \texttt{hilbert} must be present (may be empty array).
  \item If \texttt{execution.mode==="cycle"} then \texttt{hilbertPolicy} cannot be \texttt{none} (cycle mode requires trace policy \texttt{on-export} or \texttt{full-trace}).
  \item \texttt{compiled.qubitCount} is non-negative integer; \texttt{gates}/\texttt{tokenMap} present.
  \item If \texttt{deviceTable} present: shape $n\times 3$ with $n=\texttt{compiled.qubitCount}$.
  \item If \texttt{checkpoints} present: each payload's \texttt{globalState.length} $= 2^{\texttt{compiled.qubitCount}}$.
  \item If \texttt{compiled.processParams}/\texttt{returnValues} present: arrays of strings matching protocol compile metadata.
\end{enumerate}

\subsection{Upgrade system implementation sketch}

\begin{lstlisting}[style=code]
// data/formats/qpucirMigrate.ts (proposed)
export function parseQpucirPayloadV2(contents: string): QpucirPayloadV2 {
  const parsed = JSON.parse(contents);
  if (parsed.format !== 'qpucir') throw new Error('Not a qpucir payload');
  if (parsed.version === 1) return migrateV1ToV2(parsed);
  if (parsed.version === 2) {
    validateQpucirV2(parsed);
    return parsed;
  }
  throw new Error(`Unsupported qpucir version: ${parsed.version}`);
}

function migrateV1ToV2(v1: QpucirPayloadV1): QpucirPayloadV2 {
  const compiled = compileQpuProtocol(v1.source, {});
  return {
    format: 'qpucir', version: 2, name: v1.name, source: v1.source,
    exportedAt: v1.exportedAt,
    execution: { mode: 'flat' },
    hilbertPolicy: 'none',
    compiled: { ...v1.compiled, processParams: compiled.processParams,
      returnValues: compiled.returnValues },
  };
}
\end{lstlisting}

\textbf{Migration v1$\rightarrow$v2:} set \texttt{version: 2}, \texttt{execution.mode: "flat"}, \texttt{hilbertPolicy: "none"}; recompile \texttt{compiled} from \texttt{source}. Current \texttt{parseQpucirPayload} ignores \texttt{version} and reads \texttt{source} only---upgrade path preserves that for plain imports while JSON exports gain validation.

\subsection{Catalog storage}

\textbf{v1} key: \texttt{qpu-process-catalog-v1}. \textbf{v2} key: \texttt{qpu-process-catalog-v2}. On first load, import v1 entries with \texttt{qpucirVersion: 1} defaults.

\subsection{Session and correction exports}

\texttt{qpu-session} v1 and \texttt{qpu-correction-report} v1 as previously specified; amplitudes use the same \texttt{\{re,im\}} JSON rules.

\subsection{Upgrade matrix}

\begin{center}
\begin{tabular}{@{}p{3cm}p{4.5cm}p{4.5cm}}
\toprule
Object & Reader action & Writer action \\
\midrule
\texttt{qpucir v1} & Import \texttt{source}; ignore stale \texttt{compiled} if recompiling. v1 reader on v2 file: reject or strip per policy (A/B) above. & Emit v1 only if \texttt{hilbertPolicy===none} and flat mode \\
\texttt{qpucir v2} & \texttt{validateQpucirV2} then execute & Preferred when cycle/Hilbert/checkpoints used \\
\texttt{qpuio v1} & Unchanged & Unchanged \\
Catalog v1 & Migrate on read to v2 & Write v2 only \\
Plain text & \texttt{parseQpucirPayload} today & unchanged \\
\bottomrule
\end{tabular}
\end{center}

%=============================================================================
\section{Verification}
%=============================================================================

\subsection{Cross-implementation tests}

For each bundled process $P$ and truth-table row $\mathbf{x}$:
\begin{equation}
  \mathrm{Measure}(\mathrm{Run}_{\mathrm{TS}}(P,\mathbf{x})) =
  \mathrm{Measure}(\mathrm{Run}_{\mathrm{Py}}(P,\mathbf{x})).
\end{equation}

\textbf{Hilbert tests:} In cycle mode, compare sorted $(\mathit{register},c)$ keys of $\HS$ against the Python CLI memory snapshot for $c\le c_{\max}$ with \textbf{exact key-set equality}. Compare amplitude vectors element-wise: $|a_i-a'_i|\le\varepsilon_a$ for all components (using $\varepsilon_a=10^{-6}$ from \S3.1). Ket strings are informational; amplitudes are authoritative.

\subsection{Regression suite}

\begin{itemize}
  \item \texttt{hilbertSpace.test.ts} --- \texttt{output}, destructive \texttt{input}, \texttt{prune}, \texttt{view}.
  \item \texttt{cycleMode.test.ts}, \texttt{checkpoint.test.ts}, \texttt{formatQubitState.test.ts}, \texttt{exportMigration.test.ts}.
  \item \texttt{blochQuadrature.test.ts} (existing) --- fast vs grid integral.
\end{itemize}

%=============================================================================
\section{Implementation phases}
%=============================================================================

\begin{enumerate}[label=\textbf{Phase \Roman*.}]
  \item \texttt{HilbertSpace} + \texttt{formatQubitState} + cycle executor; lower \texttt{SAVE\_STATE}/\texttt{LOAD\_STATE}.
  \item \texttt{RegisterBank}, noise, device table; Module Lab Hilbert context; extend particle math kernels.
  \item \texttt{qpucir v2}, \texttt{qpucirMigrate.ts}, catalog v2.
  \item HTTP sidecar + \texttt{VITE\_QPU\_API\_URL} client.
  \item Sparse \texttt{StateVector}; WebGPU (\texttt{webGpu.ts} probe exists, no backend yet).
\end{enumerate}

%=============================================================================
\section{Summary}
%=============================================================================

Parity requires: cycle-capable executor with destructive-pop $\HS$ (Python semantics); math layers extending existing marginal $\rho_k$, Euler ket, and triple-integral $\langle\rho\rangle$; strict API contracts (errors, amplitudes, session locking); versioned exports with validation; sidecar separate from static Pages. Verified gaps in current code: HilbertSpace, cycle runtime, checkpoints, noise, and sidecar are specified here but not yet in \texttt{src/}.

\end{document}
